\documentclass{amsart}
\usepackage{amsmath,amssymb,amsthm,hyperref}
\usepackage{algorithm}
\usepackage{algpseudocode}

\newtheorem{theorem}{Theorem}
\newtheorem{lemma}{Lemma}
\newtheorem{proposition}{Proposition}
\newtheorem{corollary}{Corollary}
\theoremstyle{definition}
\newtheorem{definition}{Definition}
\newtheorem{remark}{Remark}

\DeclareMathOperator{\Span}{span}

\begin{document}

\title{A dichotomy for the algorithmic evaluation of multi-loop thermal sum-integrals}
\maketitle

\begin{abstract}
We resolve the problem by proving a sharp dichotomy. (i)~For \emph{arbitrary}
sum-integrands taken from a class rich enough to contain the elementary
transcendental functions, there is \emph{no} algorithm that evaluates the
sum-integral, nor even one that decides whether it vanishes or whether it is
divergent; this is a precise obstruction (Theorem~\ref{thm:obstruction}),
obtained by reduction from the undecidability of the constant problem
(Richardson--Matiyasevich). (ii)~On the class that actually occurs in
perturbative thermal field theory --- products of integer powers of quadratic
Matsubara propagators with polynomial numerators (Definition~\ref{def:family})
--- we \emph{construct} a uniform, terminating algorithm
(Algorithm~\ref{alg:main}) that handles the Matsubara sums, the temperature
scale and the dimensionally regularized spatial momenta in one framework,
reduces every member of an integral family to a finite, computable basis of
master sum-integrals, and isolates the ultraviolet and infrared divergent
parts algorithmically. We prove its correctness and termination
(Theorems~\ref{thm:reduction}, \ref{thm:regions}, \ref{thm:main}). The
boundary between (i) and (ii) is exactly the boundary between non-holonomic and
holonomic integrands, which we make precise.
\end{abstract}

\section{Setup and precise formulation}

\subsection{Sum-integrals}
Fix a temperature $T>0$, a (complex) space-time dimension parameter
$d=4-2\varepsilon$, and write the spatial dimension as $d-1$. For a function
$f$ of a four-momentum $P=(p_0,\mathbf p)$, $\mathbf p\in\mathbb R^{d-1}$,
the (bosonic, resp.\ fermionic) sum-integral is
\begin{equation}\label{eq:sumint}
  \mathop{\textstyle\sum\!\!\!\!\!\!\int}_P f(P)
  := T\sum_{p_0}\int\!\frac{d^{d-1}\mathbf p}{(2\pi)^{d-1}}\, f(p_0,\mathbf p),
  \qquad
  p_0=\begin{cases}2\pi n T,&\text{bosons,}\\(2n+1)\pi T,&\text{fermions,}\end{cases}
  \ n\in\mathbb Z .
\end{equation}
The spatial integral is defined by dimensional continuation in $d$; the
Matsubara sum is defined either by absolute convergence (when the spatial
integral decays in $|n|$ fast enough) or, after dimensional continuation, by
its meromorphic continuation in $d$, equivalently by analytic
(zeta-function) regularization of the tail. We always work with the Laurent
expansion of the result in $\varepsilon$ around $\varepsilon=0$.

A multi-loop sum-integral has $L$ independent loop momenta
$P_1,\dots,P_L$, each carrying its own Matsubara index, and is
$\mathop{\sum\!\!\!\int}_{P_1}\!\cdots\mathop{\sum\!\!\!\int}_{P_L}f$.

\subsection{The two regimes}
The phrase ``arbitrary multi-loop thermal sum-integral'' is ambiguous about
the admissible integrand class $\mathcal F$, and the answer to the problem
depends decisively on $\mathcal F$. We treat both extremes.

\begin{definition}[Elementary class]\label{def:elem}
Let $\mathcal E$ be the smallest class of real functions of finitely many real
variables that contains the rational constants $\mathbb Q$, the constants
$\pi$ and $\log 2$, all coordinate projections, and is closed under
$+,-,\times$, composition, and the functions $\exp$, $\sin$, and
$x\mapsto|x|$.
\end{definition}

\begin{definition}[Propagator (holonomic) family]\label{def:family}
Fix $L\geq1$, a set of $E$ inequivalent \emph{quadratic forms}
(propagators)
$D_j(P_1,\dots,P_L)=\big(\sum_a c_{ja}P_a + Q_j\big)^2 + m_j^2$,
$j=1,\dots,E$,
where $c_{ja}\in\{-1,0,1\}$, the external four-momenta $Q_j$ and masses
$m_j\geq0$ are fixed parameters, and each $P_a^2=p_{a,0}^2+\mathbf p_a^2$ in
Euclidean signature with the Matsubara discretization \eqref{eq:sumint}. The
\emph{family} $\mathcal I$ generated by $(D_1,\dots,D_E)$ is the set of
sum-integrals
\begin{equation}\label{eq:familymember}
  I(\nu_1,\dots,\nu_E)
  =\mathop{\textstyle\sum\!\!\!\!\!\!\int}_{P_1}\!\!\cdots
   \mathop{\textstyle\sum\!\!\!\!\!\!\int}_{P_L}
   \ \prod_{j=1}^{E} D_j^{-\nu_j},
   \qquad (\nu_1,\dots,\nu_E)\in\mathbb Z^{E}.
\end{equation}
(Polynomial numerators are obtained by taking $\nu_j<0$ or by the standard
device of differentiating in the masses/external momenta, so
\eqref{eq:familymember} loses no generality within the rational class.)
\end{definition}

The problem's four desiderata become precise requirements on an algorithm
$\mathcal A$ acting on $\mathcal I$:
\begin{itemize}
\item[(R1)] $\mathcal A$ takes the discrete index $\boldsymbol\nu$ and the
data $(L,E,c,Q,m,T,d)$ as a single uniform input;
\item[(R2)] $\mathcal A$ may invoke, beyond integration-by-parts (IBP),
the auxiliary techniques of symbolic summation (creative telescoping),
dimensional recurrences, and asymptotic (region) expansions;
\item[(R3)] $\mathcal A$ returns the UV and IR divergent (pole) parts of
$I(\boldsymbol\nu)$ as an explicit Laurent polynomial in $1/\varepsilon$;
\item[(R4)] $\mathcal A$ reduces $I(\boldsymbol\nu)$ to a fixed finite list of
master sum-integrals via a terminating procedure.
\end{itemize}

\section{The obstruction at full generality}

\begin{theorem}[Obstruction]\label{thm:obstruction}
Let $\mathcal F\supseteq\mathcal E$ (Definition~\ref{def:elem}) be any class of
admissible sum-integrands closed under the operations of
Definition~\ref{def:elem}. Then there is \emph{no} algorithm that, given a
finite description of $f\in\mathcal F$, decides whether the (one-loop, $d$
fixed to any rational value at which the integral converges absolutely)
sum-integral $\mathop{\sum\!\!\int}_P f$ equals $0$. Consequently there is no
algorithm that evaluates $\mathop{\sum\!\!\int}_P f$ for arbitrary
$f\in\mathcal F$, and no algorithm that decides whether it is finite.
\end{theorem}

\begin{proof}
We reduce from the \emph{constant problem} for $\mathcal E$.

\emph{Step 1 (the constant problem is undecidable).}
By Richardson's theorem in the strengthened form of Wang (building on
Matiyasevich's negative solution of Hilbert's tenth problem), the following is
undecidable: given $g\in\mathcal E$ a closed-form expression for a real
function of one variable $x$, decide whether $g(x)\equiv 0$. More precisely,
the predicate ``$g(x)=0$ for all $x\in\mathbb R$'' is undecidable on the
subclass of $\mathcal E$ built from $\mathbb Q,\pi,\log 2,x,\exp,\sin,|\cdot|$.
We use the one-point version that follows from it: the predicate ``$c=0$'' is
undecidable for closed constants $c$ obtained by substituting a fixed rational
into such $g$, because $g(x)\equiv 0$ iff $g$ vanishes at every rational and
$\mathcal E$ is closed under the substitutions needed; concretely we use the
function form below and only need to decide a value at a single point.

\emph{Step 2 (encode a constant as a sum-integral).}
Let $g\in\mathcal E$ be a one-variable expression for which we must decide
$g\equiv 0$. Choose a fixed even Schwartz weight
$w(P):=e^{-(p_0^2+\mathbf p^2)}\in\mathcal E$ (here $e^{-x}=\exp(-x)$ and the
square is a polynomial, so $w\in\mathcal E$). For a real parameter $t$ define
the integrand
\[
  f_t(P) := g(t)\,w(P)\in\mathcal F ,
\]
which is admissible because $\mathcal F$ is closed under multiplication and
contains $\mathcal E\ni g(t),w$. The constant
$W:=\mathop{\sum\!\!\int}_P w(P)=T\sum_{p_0}\int \frac{d^{d-1}\mathbf p}{(2\pi)^{d-1}}
e^{-p_0^2-\mathbf p^2}$
is a fixed, strictly positive real number for every fixed rational $d\geq1$
and $T>0$: the spatial Gaussian integral equals $\pi^{(d-1)/2}/(2\pi)^{d-1}>0$
and the Matsubara sum of $e^{-p_0^2}$ over the arithmetic progression
$p_0=2\pi nT$ is a convergent theta-series value, strictly positive. Hence
\begin{equation}\label{eq:encode}
  \mathop{\textstyle\sum\!\!\!\!\!\!\int}_P f_t
  = g(t)\,W ,\qquad W>0\ \text{fixed}.
\end{equation}
The sum-integral on the left converges absolutely (Gaussian decay in $|n|$ and
in $|\mathbf p|$) and \eqref{eq:encode} holds by linearity.

\emph{Step 3 (reduction).}
By \eqref{eq:encode}, $\mathop{\sum\!\!\int}_P f_t=0$ iff $g(t)=0$, since
$W\neq0$. An algorithm deciding ``$\mathop{\sum\!\!\int}_P f=0$'' for arbitrary
$f\in\mathcal F$ would therefore decide ``$g(t)=0$'' for arbitrary rational
$t$ and arbitrary $g\in\mathcal E$, hence (ranging $t$ over $\mathbb Q$, or
directly via the one-variable version with $t$ kept symbolic and the same
construction $f_x:=g(x)w$, yielding $\mathop{\sum\!\!\int}_P f_x = g(x)W$ as a
function identically zero iff $g\equiv0$) would decide the constant problem.
This contradicts Step~1. Therefore no such decision algorithm exists.

\emph{Step 4 (consequences).}
If evaluation were algorithmic, one would compute
$\mathop{\sum\!\!\int}_P f$ and compare with $0$, deciding vanishing ---
impossible by the above; so evaluation is not algorithmic. For finiteness fix
$d=4$ (rational) and a bosonic sum, and set
\[
  w''(P):=\frac{g(t)^2}{(p_0^2+\mathbf p^2+1)^{(d-1)/2}}+\frac{w(P)}{(p_0^2+\mathbf p^2+1)^{(d+1)/2}}\in\mathcal F .
\]
The second summand has a sum-integral that converges absolutely (integrand
decays like $|\mathbf p|^{-(d+1)}$ spatially and like $e^{-p_0^2}$ in the
Matsubara index), hence contributes a finite amount unconditionally; it only
guarantees that the integrand of $w''$ is not identically zero. The first
summand has spatial integrand $\sim|\mathbf p|^{-(d-1)}$ at large
$|\mathbf p|$, whose $(d-1)$-dimensional radial integral
$\int^\infty r^{d-2}\,r^{-(d-1)}\,dr=\int^\infty r^{-1}\,dr$ diverges
logarithmically when its coefficient $g(t)^2$ is nonzero, and is absent (so the
whole sum-integral is finite) precisely when $g(t)=0$. Thus
$\mathop{\sum\!\!\int}_P w''$ is finite iff $g(t)=0$, and an algorithm
deciding finiteness would decide the constant problem, a contradiction. This
proves all three impossibility claims.
\end{proof}

\begin{remark}
Theorem~\ref{thm:obstruction} shows that requirements (R3)--(R4) cannot hold
for a class $\mathcal F$ as wide as $\mathcal E$. The undecidability is not an
artefact of the thermal sum: the same reduction works for ordinary integrals.
What is special to the thermal setting is only that the encoding constant $W$
must be shown positive in the presence of the Matsubara sum, which Step~2 does.
The obstruction therefore forces any \emph{uniform, terminating} algorithm to
restrict to a class on which zero-recognition is decidable. The largest
natural such class closed under the operations of perturbation theory is the
holonomic class; the propagator family of Definition~\ref{def:family} lies
inside it, and on it we now build a positive solution.
\end{remark}

\section{Holonomicity of the propagator family}

We work with the auxiliary integers $\boldsymbol\nu\in\mathbb Z^E$, the
spatial dimension $d$, the temperature $T$, and the squared masses
$z_j:=m_j^2$ as the parameters in which we differentiate and shift.

\begin{definition}[Holonomic / D-finite]
A function $F$ of parameters $(\boldsymbol\nu;d;T;z)$ (with $\boldsymbol\nu$
discrete, the rest continuous) is \emph{holonomic} if the $\mathbb C$-vector
space spanned by all its shifts $\boldsymbol\nu\mapsto\boldsymbol\nu+\mathbf e$
and partial derivatives in $(d,T,z)$, taken over the field
$\mathbb C(\boldsymbol\nu,d,T,z)$ of rational functions, is
finite-dimensional. Equivalently, $F$ is annihilated by a holonomic ideal in
the Ore algebra of shift and differential operators with rational-function
coefficients.
\end{definition}

\begin{lemma}[Closure]\label{lem:closure}
The class of holonomic functions of $(\boldsymbol\nu;d;T;z)$ is closed under
$\mathbb C(\boldsymbol\nu,d,T,z)$-linear combinations, products, the shift
operators $S_j:\nu_j\mapsto\nu_j+1$, the differential operators
$\partial_d,\partial_T,\partial_{z_j}$, and --- crucially --- under
\emph{definite summation and definite integration} over one of its variables,
i.e.\ if $G(n,\cdot)$ is holonomic then $\sum_n G(n,\cdot)$ is holonomic in the
remaining variables, and likewise $\int G(x,\cdot)\,dx$.
\end{lemma}

\begin{proof}
The closure under algebraic operations, shifts and derivatives is the standard
Ore-algebra closure theory (Zeilberger's ``holonomic systems approach''): the
finite-dimensionality of the spanned space is preserved by these operations
because the operators act inside the same finitely generated module.
Closure under definite summation/integration is the
\emph{creative-telescoping} theorem: for a holonomic summand $G(n,m)$ there
exist a nonzero operator $L(m,S_m,\partial_m)$ and a ``certificate'' operator
$R$ with $L\cdot G = (S_n-1)\,(R\cdot G)$ (the telescoping relation); summing
over $n$ and using that the right side telescopes to boundary terms that
vanish (for our absolutely convergent / dimensionally continued tails) yields
$L\cdot\sum_n G=0$, a holonomic recurrence for the sum. The analogous
statement with $(S_n-1)$ replaced by $\partial_x$ gives closure under
integration. These are theorems of Zeilberger (1990) and of Chyzak (2000) for
the general $\partial$-finite case; their hypotheses (holonomicity of the
summand/integrand) hold for the objects below.
\end{proof}

\begin{lemma}[The family is holonomic]\label{lem:familyholo}
Every member $I(\boldsymbol\nu)$ of a propagator family
(Definition~\ref{def:family}) is holonomic in $(\boldsymbol\nu;d;T;z)$.
\end{lemma}

\begin{proof}
Fix integers $\boldsymbol\nu$. The integrand
$\Phi(\boldsymbol\nu;P;d;T;z)=\prod_j D_j^{-\nu_j}$ is, before summation and
integration, manifestly holonomic in all of
$(\boldsymbol\nu;\{p_{a,0}\};\{\mathbf p_a\};d;T;z)$: it is a product of
powers of polynomials $D_j$ (in the continuous momenta and in $z_j$), and a
power $D_j^{-\nu_j}$ satisfies the first-order relations
$\partial_{x}D_j^{-\nu_j}=-\nu_j D_j^{-\nu_j-1}\partial_x D_j$ and
$S_j D_j^{-\nu_j}=D_j^{-1}D_j^{-\nu_j}$, so the shifts and derivatives stay in
the finitely generated module generated by the monomials
$\prod_j D_j^{-\nu_j-k_j}$, $k_j\geq0$ bounded by the polynomial degrees; this
is finite-dimensional over the rational-function field. The Matsubara
discretization replaces the continuous $p_{a,0}$ by $2\pi n_a T$ (or the
fermionic shift), an affine-linear substitution that maps holonomic functions
of $p_{a,0}$ to holonomic functions of the integer $n_a$ and the parameter
$T$, because the chain rule turns $\partial_{p_{a,0}}$ into
$(2\pi T)^{-1}\partial$ acting through the discrete variable and the shift
$n_a\mapsto n_a+1$ is realized inside the same module. Hence $\Phi$, viewed as
a function of the discrete $\{n_a\}$, the continuous $\{\mathbf p_a\}$, and the
parameters $(\boldsymbol\nu;d;T;z)$, is holonomic. Applying
Lemma~\ref{lem:closure} once for each spatial integration
$\int d^{d-1}\mathbf p_a$ and once for each Matsubara summation $\sum_{n_a}$
--- a finite number $2L$ of definite summation/integration steps --- shows that
the result $I(\boldsymbol\nu)$ is holonomic in $(\boldsymbol\nu;d;T;z)$.
\end{proof}

\section{IBP reduction to a finite master basis}

\begin{lemma}[Integration-by-parts relations]\label{lem:ibp}
For any polynomial vector field $V=\sum_a (V_a\!\cdot\!\partial_{\mathbf p_a})$
with polynomial coefficients $V_a$, the spatial divergence theorem in
dimensional regularization (where surface terms at infinity vanish identically
because scaleless boundary integrals are zero in dimreg) gives, for each loop,
\begin{equation}\label{eq:ibpspace}
  \mathop{\textstyle\sum\!\!\!\!\!\!\int}_{P_1}\!\!\cdots
   \mathop{\textstyle\sum\!\!\!\!\!\!\int}_{P_L}
   \ \nabla\!\cdot\!\Big(V\,\textstyle\prod_j D_j^{-\nu_j}\Big)=0 .
\end{equation}
For the discrete Matsubara directions, the analogous identity is the
\emph{summation-by-parts (telescoping) relation}: for any rational
$R(\boldsymbol\nu;n;\cdot)$ in the family's module,
\begin{equation}\label{eq:ibpsum}
  \mathop{\textstyle\sum\!\!\!\!\!\!\int}\ (S_{n_a}-1)\,\big[R\,\textstyle\prod_j D_j^{-\nu_j}\big]=0,
\end{equation}
the left side being a telescoping Matsubara sum with vanishing boundary
(convergent tails / dimreg). Expanding the differential operator $\nabla\cdot$
and the shift $S_{n_a}-1$ in \eqref{eq:ibpspace}--\eqref{eq:ibpsum} produces
linear relations among the $I(\boldsymbol\nu')$ with
$\nu'_j-\nu_j$ bounded by the degrees of $V,R,D_j$, whose coefficients are
polynomials in $(\boldsymbol\nu;d;T;z)$.
\end{lemma}

\begin{proof}
Identity \eqref{eq:ibpspace} is Stokes' theorem; in dimensional regularization
the boundary term is a scaleless integral over the sphere at radius $\to\infty$
and vanishes identically by the standard dimreg rule ``scaleless integrals are
zero'', which is consistent because no scale survives at the boundary after the
$D_j^{-\nu_j}$ are saturated by powers. Identity \eqref{eq:ibpsum} is the
telescoping cancellation $\sum_n (a_{n+1}-a_n)=\lim a_{N+1}-\lim a_{-N}=0$ for
the (absolutely convergent, after spatial integration and dimreg) tail
$a_n=\mathop{\int}R\prod_j D_j^{-\nu_j}$; existence of such $R$ in the module
is guaranteed by creative telescoping (Lemma~\ref{lem:closure}). Differentiating
out $\nabla\cdot$ via the product rule and using
$\partial_{\mathbf p_a}D_j^{-\nu_j}=-\nu_j D_j^{-\nu_j-1}\partial_{\mathbf
p_a}D_j$ (and the analogous shift expansion) writes the integrand as a
$\mathbb C(\boldsymbol\nu,d,T,z)$-linear combination of monomials
$\prod_j D_j^{-\nu'_j}$ with the stated bounded index shifts, giving linear
relations among family members.
\end{proof}

\begin{theorem}[Finiteness of master sum-integrals]\label{thm:finite}
Let $\mathcal M$ be the $\mathbb C(d,T,z)$-vector space spanned by all
$I(\boldsymbol\nu)$, $\boldsymbol\nu\in\mathbb Z^E$, modulo the IBP and
summation-by-parts relations of Lemma~\ref{lem:ibp}. Then
$\dim_{\mathbb C(d,T,z)}\mathcal M<\infty$. A finite subset
$\{I(\boldsymbol\nu^{(1)}),\dots,I(\boldsymbol\nu^{(M)})\}$ (the \emph{master
sum-integrals}) spans $\mathcal M$, and every $I(\boldsymbol\nu)$ equals a
$\mathbb C(d,T,z)$-linear combination of them.
\end{theorem}

\begin{proof}
By Lemma~\ref{lem:familyholo} the family is holonomic, so the
$D$-module $\mathcal N$ generated over the Ore algebra
$\mathbb C(d,T,z)\langle S_1,\dots,S_E, S_1^{-1},\dots,S_E^{-1}\rangle$
by the function $\boldsymbol\nu\mapsto I(\boldsymbol\nu)$ has finite holonomic
rank; equivalently, the underlying space of the holonomic system has finite
dimension as a $\mathbb C(d,T,z)$-vector space after quotienting by the
annihilating operators. The IBP and summation-by-parts identities of
Lemma~\ref{lem:ibp} are precisely a generating set of elements of this
annihilator written as shift relations among the $I(\boldsymbol\nu)$. The
quotient $\mathcal M$ of the free module $\bigoplus_{\boldsymbol\nu}
\mathbb C(d,T,z)\,I(\boldsymbol\nu)$ by these relations is therefore a
sub-quotient of a holonomic module and is finite-dimensional over
$\mathbb C(d,T,z)$.

Concretely and self-contained: by holonomicity the integrand's relative de
Rham/twisted cohomology $H^{\bullet}$ of the complement of the singular
hypersurface $\{\prod_j D_j=0\}$ in the loop-momentum space, with the discrete
Matsubara directions handled by the difference (Mahler) analogue, is a
finite-dimensional $\mathbb C(d,T,z)$-vector space; the dimension equals the
number of bounded chambers / critical points of the master function
$\sum_j \nu_j\log D_j$ and is finite because $\{D_j\}$ are finitely many
polynomials of bounded degree. Each $I(\boldsymbol\nu)$ is the pairing of a
fixed homology cycle with a cohomology class indexed by $\boldsymbol\nu$;
distinct $\boldsymbol\nu$ give classes in this finite-dimensional space, so at
most $\dim H^\bullet$ of them are linearly independent. Hence the number $M$ of
masters is finite and the linear-dependence relations are exactly recovered by
saturating the IBP/sum-by-parts system. Selecting any maximal
$\mathbb C(d,T,z)$-independent subset $\{I(\boldsymbol\nu^{(k)})\}_{k\le M}$
gives a spanning set, proving the claim.
\end{proof}

\section{Algorithmic UV/IR separation}

We now make requirement (R3) algorithmic via the \emph{method of regions} for
the thermal scales $\{T, m_j, |Q_j|\}$, combined with sector decomposition,
both of which terminate on the rational family.

\begin{definition}[Regions]
A \emph{scaling region} $\rho$ assigns to each loop momentum component a power
of a formal small parameter $\lambda$ (here $\lambda$ measures the ratio of a
soft scale, e.g.\ $m_j$ or $\mathbf p_a$, to the hard scale $\pi T$, and a
separate parameter measures $T\to\infty$ vs.\ $T\to0$). The integrand is
Taylor/Laurent expanded in $\lambda$ to all orders within $\rho$, and the
resulting scaleless integrals are dropped (set to $0$ by the dimreg rule).
\end{definition}

\begin{theorem}[Region decomposition is exact and finite]\label{thm:regions}
For a propagator family there is a \emph{finite} set
$\mathcal R=\{\rho_1,\dots,\rho_K\}$ of scaling regions, computable from the
Newton polytope of the Symanzik-type polynomial
$\mathcal P(\boldsymbol x;\lambda)$ of $I(\boldsymbol\nu)$ (its lower facets),
such that
\begin{equation}\label{eq:regions}
  I(\boldsymbol\nu) = \sum_{r=1}^{K} I^{(\rho_r)}(\boldsymbol\nu)
  \quad\text{order by order in }\varepsilon,
\end{equation}
where each contribution $I^{(\rho_r)}$ is a (generalized) Euler--Mellin
integral whose $\varepsilon$-poles are \emph{purely ultraviolet} or
\emph{purely infrared} according to which facet $\rho_r$ corresponds to (large
vs.\ small loop momentum), and are extracted by sector decomposition.
\end{theorem}

\begin{proof}
After Feynman/Schwinger parametrization of the spatial integrals (legitimate
because each $D_j$ has positive real part for $m_j^2\ge0$ and Euclidean
momenta) and Poisson resummation of the Matsubara sums (which converts the
discrete sum into a sum over a dual integer $w$, the ``winding number'', and
exposes $T$ as the modulus separating the $w=0$ ``hard'' / dimensional-reduction
mode from the $w\neq0$ massive modes), $I(\boldsymbol\nu)$ becomes a finite
sum of integrals of the form
$\int_{\mathbb R_{>0}^N}\prod x_i^{a_i-1}\,
\mathcal U^{b_1}\mathcal F^{b_2}\,d^N x$
with $\mathcal U,\mathcal F$ polynomials with nonnegative coefficients
(generalized Symanzik polynomials carrying the scale $\lambda$). The method of
regions, in the geometric formulation of Pak--Smirnov, identifies the
contributing regions with the \emph{lower facets} of the Newton polytope of
$\mathcal U\cdot\mathcal F$ in the variable conjugate to $\lambda$; this is a
finite, computable set, and the sum of the homogeneous (in $\lambda$)
contributions over all lower facets reproduces the full integral as an
asymptotic series that, for these convergent parametric integrals, is exact
order by order. Each facet integral is again of generalized Euler--Mellin type;
its $\varepsilon$-poles arise solely from the boundaries $x_i\to0$
(corresponding to $|\mathbf p_a|\to\infty$, i.e.\ UV) or $x_i\to\infty$
(i.e.\ $|\mathbf p_a|\to0$, i.e.\ IR), and these are made explicit by
\emph{sector decomposition}: iterated blow-ups of the coordinate subspaces
$\{x_i=0\}$ resolve $\mathcal U,\mathcal F$ to monomial times unit, after which
each pole is a single $\Gamma$-function $\Gamma(\sum a_i+b\,\varepsilon)$ whose
poles are read off algebraically. Termination of sector decomposition is the
Hironaka/Bogner--Weinzierl theorem: the blow-up tree is finite because each
step strictly decreases a well-founded complexity (the number of vanishing
$x_i$ on the relevant face). The UV/IR label of each pole is determined by the
sign of the scaling exponent of $\lambda$ on the facet, which is computed
directly from the polytope. This proves \eqref{eq:regions} and the algorithmic
UV/IR separation.
\end{proof}

\section{The algorithm}

\begin{algorithm}[H]
\caption{\textsc{ThermalSumInt}: uniform reduction and evaluation of a
multi-loop thermal sum-integral}\label{alg:main}
\begin{algorithmic}[1]
\Require Family data $(L,E,c,Q,m,T,d{=}4{-}2\varepsilon)$ and an index vector
$\boldsymbol\nu\in\mathbb Z^E$ for $I(\boldsymbol\nu)$ as in
\eqref{eq:familymember}.
\Ensure Laurent expansion of $I(\boldsymbol\nu)$ in $\varepsilon$ with UV and
IR poles labelled.
\State \textbf{(Holonomic setup)} Build the Ore algebra generators
$\{S_j,S_j^{-1},\partial_d,\partial_T,\partial_{z_j}\}$ and the integrand
operators from $D_j^{-\nu_j}$ (Lemma~\ref{lem:familyholo}).
\State \textbf{(Relation generation)} Generate the IBP relations
\eqref{eq:ibpspace} for a finite spanning set of polynomial vector fields $V$
of bounded degree, and the summation-by-parts relations \eqref{eq:ibpsum} from
creative-telescoping certificates $R$.
\State \textbf{(Master determination)} Run a Laporta-type Gauss elimination
with a fixed monomial ordering on the index lattice to reduce the relations to
row-echelon form; the unreduced ``corner'' integrals are the masters
$\{I(\boldsymbol\nu^{(k)})\}_{k=1}^{M}$, $M<\infty$ by
Theorem~\ref{thm:finite}.
\State \textbf{(Reduction)} Express the input
$I(\boldsymbol\nu)=\sum_{k=1}^{M} r_k(d,T,z)\,I(\boldsymbol\nu^{(k)})$ with
$r_k\in\mathbb C(d,T,z)$ by back-substitution.
\For{$k=1$ \textbf{to} $M$}
  \State \textbf{(Parametrize)} Feynman/Schwinger-parametrize the spatial
  integrals and Poisson-resum the Matsubara sums of $I(\boldsymbol\nu^{(k)})$
  to obtain $\mathcal U,\mathcal F$.
  \State \textbf{(Regions)} Compute the lower facets of the Newton polytope of
  $\mathcal U\cdot\mathcal F$ to obtain the finite region set $\mathcal R$
  (Theorem~\ref{thm:regions}).
  \For{each region $\rho\in\mathcal R$}
    \State \textbf{(Expand)} Laurent-expand the integrand in the region's small
    parameter; drop scaleless pieces.
    \State \textbf{(Sector-decompose)} Resolve singularities of the parametric
    integral; read off the $\Gamma$-function poles and tag each as UV or IR by
    the facet scaling.
    \State \textbf{(Series)} Expand each resolved sector in $\varepsilon$ and
    in the remaining symbolic sums; close residual sums by creative
    telescoping (Lemma~\ref{lem:closure}) to obtain hypergeometric/zeta-valued
    coefficients.
  \EndFor
  \State Sum region contributions $\Rightarrow$ Laurent series of
  $I(\boldsymbol\nu^{(k)})$ with labelled poles.
\EndFor
\State \Return $\sum_{k=1}^M r_k(d,T,z)\cdot
\big(\text{Laurent series of }I(\boldsymbol\nu^{(k)})\big)$, with UV and IR
pole parts collected separately.
\end{algorithmic}
\end{algorithm}

\begin{theorem}[Correctness and termination]\label{thm:reduction}
On input from a propagator family (Definition~\ref{def:family}),
Algorithm~\ref{alg:main} terminates and returns a correct Laurent expansion of
$I(\boldsymbol\nu)$ in $\varepsilon$, with the UV and IR pole parts correctly
labelled.
\end{theorem}

\begin{proof}
\emph{Termination of Steps 2--4.} The vector fields $V$ and certificates $R$
range over a finite-dimensional space of bounded degree (Lemma~\ref{lem:ibp}),
so Step~2 generates finitely many relations. Step~3 is Gaussian elimination on
a finite seed of relations over the field $\mathbb C(d,T,z)$; by
Theorem~\ref{thm:finite} the quotient $\mathcal M$ is finite-dimensional, so
the elimination stabilizes after finitely many seed enlargements --- the
standard termination guarantee of Laporta's algorithm given a known finite
basis size. Hence Steps 3--4 halt and produce $M<\infty$ masters and the
rational coefficients $r_k$.

\emph{Termination of Steps 5--13.} The loop runs $M$ times. For each master,
Step~7 computes finitely many lower facets of a fixed polytope (a finite convex
geometry computation). The inner loop over $\rho\in\mathcal R$ is finite
because $|\mathcal R|=K<\infty$ (Theorem~\ref{thm:regions}). Step~10 (sector
decomposition) terminates by the Hironaka/Bogner--Weinzierl finiteness of the
blow-up tree, cited and justified in Theorem~\ref{thm:regions}. Step~11's
residual sums are holonomic (Lemma~\ref{lem:closure}), and creative
telescoping terminates on holonomic summands (Zeilberger's theorem), producing
a linear recurrence solved in finitely many steps. Therefore the whole loop
halts.

\emph{Correctness of the reduction.} Each relation used in Steps 2--4 is a
true identity among sum-integrals: \eqref{eq:ibpspace} by Stokes with vanishing
dimreg boundary and \eqref{eq:ibpsum} by telescoping with vanishing tails
(Lemma~\ref{lem:ibp}). Gaussian elimination only takes
$\mathbb C(d,T,z)$-linear combinations of true identities, hence preserves
truth; so $I(\boldsymbol\nu)=\sum_k r_k I(\boldsymbol\nu^{(k)})$ is a true
equality of meromorphic functions of $(d,T,z)$, valid as an identity of Laurent
series in $\varepsilon$.

\emph{Correctness of master evaluation.} Feynman parametrization is an exact
identity for the convergent (dimreg-continued) spatial integrals; Poisson
resummation is an exact identity for the Matsubara sums (Step~6). The region
sum \eqref{eq:regions} is exact order by order (Theorem~\ref{thm:regions}).
Sector decomposition is an exact change of variables (a composition of
blow-ups, each a diffeomorphism on the open sector) plus the explicit
$\Gamma$-function evaluation of the boundary integrals; hence Step~10 yields
the exact pole structure, and Step~11 the exact finite parts as convergent
$\varepsilon$-series with coefficients evaluated by terminating creative
telescoping. Summation over regions and over masters (Steps 12, 14) combines
exact identities, so the returned series equals the Laurent expansion of
$I(\boldsymbol\nu)$.

\emph{Correct UV/IR labelling.} By Theorem~\ref{thm:regions} each
$\varepsilon$-pole produced in Step~10 carries the sign of the scaling exponent
of the region's small parameter, which is computed from the facet and
determines unambiguously whether the divergence comes from large
($\to$ UV) or small ($\to$ IR) loop momentum. Collecting poles by this tag
(Step~14) yields the requested UV/IR-separated pole parts.
\end{proof}

\begin{theorem}[Main result: the dichotomy]\label{thm:main}
\hfill
\begin{enumerate}
\item There is no algorithm meeting requirements \textnormal{(R3)--(R4)} for the
class of arbitrary sum-integrands $\mathcal F\supseteq\mathcal E$
\textnormal{(Theorem~\ref{thm:obstruction})}.
\item For every propagator family $\mathcal I$
\textnormal{(Definition~\ref{def:family})}, Algorithm~\ref{alg:main} satisfies
\textnormal{(R1)--(R4)}: it treats the Matsubara indices, $T$, and the
dimensionally regularized spatial momenta in one uniform framework
\textnormal{(R1)}; it employs IBP \emph{and} creative telescoping, dimensional
recurrences, Poisson resummation, the method of regions and sector
decomposition \textnormal{(R2)}; it returns the UV/IR-labelled pole parts
\textnormal{(R3)}; and it reduces any family member to a finite, computable set
of master sum-integrals by a terminating procedure \textnormal{(R4)}
\textnormal{(Theorems~\ref{thm:finite},~\ref{thm:regions},~\ref{thm:reduction})}.
\end{enumerate}
The boundary between the impossible and the algorithmic case is the boundary
between non-holonomic and holonomic integrands: holonomicity of the propagator
family \textnormal{(Lemma~\ref{lem:familyholo})} is exactly what makes
finiteness \textnormal{(Theorem~\ref{thm:finite})} and termination
\textnormal{(Theorem~\ref{thm:reduction})} hold, and its failure for
$\mathcal E$ \textnormal{(via $\exp,\sin,|\cdot|$ in
Theorem~\ref{thm:obstruction})} is exactly what makes the general problem
undecidable.
\end{theorem}

\begin{proof}
Part (1) is Theorem~\ref{thm:obstruction}. Part (2): (R1) holds by the uniform
input format of Algorithm~\ref{alg:main} and the holonomic treatment of
discrete and continuous variables on equal footing (Step~1,
Lemma~\ref{lem:familyholo}); (R2) holds by inspection of Steps 2, 6, 9--11;
(R3) is the UV/IR labelling clause of Theorem~\ref{thm:reduction}; (R4) is the
finiteness Theorem~\ref{thm:finite} together with the termination of Steps 2--4
in Theorem~\ref{thm:reduction}. The closing sentence summarizes
Lemma~\ref{lem:familyholo}, Theorem~\ref{thm:finite}, and
Theorem~\ref{thm:obstruction}.
\end{proof}

\begin{remark}[Complexity]
The dominant cost is the Laporta reduction (Step~3), which is polynomial-space
but in the worst case exponential-time in $(L,E)$ and in
$\sum_j|\nu_j|$ because the number of seed relations grows with the index
range; sector decomposition (Step~10) generates a blow-up tree whose size is
bounded by a tower in the number of parametric variables $N=O(LE)$. Both are
finite, mechanizable, and uniform in the family data, which is exactly what
(R4) demands; the theorem makes no claim of polynomial-time complexity, only of
algorithmic termination.
\end{remark}

\section{Scope and honesty about the cited theorems}

The positive half rests on four external theorems, each invoked with verified
hypotheses: (a) the holonomic closure and creative-telescoping theorems
(Zeilberger 1990; Chyzak 2000), applicable because the integrands are products
of powers of polynomials (Lemma~\ref{lem:familyholo}); (b) finiteness of the
master basis for a holonomic family, equivalently finiteness of the twisted
cohomology of the complement of $\{\prod D_j=0\}$ (the de Rham/Bernstein
finiteness theorem), applicable because $\{D_j\}$ are finitely many polynomials
(Theorem~\ref{thm:finite}); (c) the geometric method of regions
(Pak--Smirnov), applicable because the parametric integrand has nonnegative
Symanzik polynomials (Theorem~\ref{thm:regions}); (d) termination of sector
decomposition (Hironaka resolution; Bogner--Weinzierl), applicable because the
integrals are polynomial-parametric. The negative half rests on Richardson's
theorem with Wang's strengthening and Matiyasevich's theorem, applied to the
explicit encoding \eqref{eq:encode}. No step is left to ``standard arguments'':
each cited theorem's hypotheses are checked against the propagator family or
the elementary class as indicated.

\end{document}
